\chapter{前言}

真传一句话，假传万卷书。具体到计算机里的某一项技能，一句话当然说不完。但一个人想学计算机体系结构，也别上来就给他推荐《计算机体系结构：量化研究方法》（\emph{Computer Architecture: A Quantitative Approach}），这谁能看懂。

我想先讲讲自己是怎么走到 CPU 这里来的。

小时候，我很想知道计算机是怎么工作的，但从来不敢去想硬件是怎么造出来的。在我眼里，硬件像是凭空长出来的魔法，就像这个世界本来的设定一样摆在那里。我不会去质疑它，甚至不会去好奇它是怎么工作的。我只知道 CPU 一年比一年快，自己的电脑很快就落伍了。

到了高中和大学，我开始学黑客技术。对我来说，它是一条通往底层软件的通道，是我的“白兔”——就像《爱丽丝梦游仙境》里那只白兔，跟着它，就掉进了另一个世界。这一次，我终于敢去好奇了：那些黑客工具是怎么做出来的？

我先后学了 Pascal、BASIC，后来是 C。我知道了“API”这个词，它是操作系统提供的服务；顺便也才弄明白，原来还有一个东西叫操作系统。等到我能调用 Win32 API 的时候，比起在 C 里用 \texttt{printf} 打印三角形，我觉得自己一下子有了力量，好像在计算机的世界里，没有什么是我做不到的。

然后是汇编，是操作系统内核的代码。我读泄露出来的 Windows 内核代码，也读开源的 Linux 内核代码。我试着写驱动，也想过写一个小小的操作系统。这件事我其实没有真的做成，但我学过赵炯的《Linux 内核完全注释》。这本书逐行注释了 Linux 0.11 内核的源代码。大学时，我还在一台真实的 PC 上装起了 Linux 0.11 内核。

我还自己写过一个 TCP 协议栈。当时我写的 rootkit 能绕过所有商业软件防火墙，但传数据很慢，因为我的做法是发一个包、确认一次，再发下一个。为了解决这个问题，我从 4.4BSD 的代码里学会了滑动窗口。

做过不少 rootkit 项目之后，我的目光转向了更底层的东西，比如 CPU 是怎么工作的。x86 汇编只给了我一点线索，离真正弄懂 CPU 还差得很远：它是怎么工作的，长什么样子，我连想象都想象不出来。

那时我参加了一个项目：用软件模拟一颗 x86 CPU，在里面运行恶意代码。我负责内存模拟的部分。不过这部分做的大多还是操作系统层面的事，比如分页、保护页（guard page）、写时复制（copy-on-write）。因为内存其实不需要真的去模拟——模拟器本身就运行在内存里。但这正是我需要的那条线索。古人说“阅藏知津”，读经要先知道渡口在哪里。这个项目，就是我通往大海的那个渡口。

又过了几年，我自己写了一个 x86 CPU 模拟器。它并不完整，但我从中弄明白了怎样解析二进制代码，一条指令是怎样由操作码、ModR/M 字节这些部分组成的。我的 CPU 之旅，就是从这里真正开始的。

我不想让这些话听起来像是我已经走到了终点，好像已经知道怎么造出一颗现代 CPU。我离那里还很远。但我确实在往那里走，而且这条路在我眼前越来越清楚。

这本书，就是我沿着这条路留下的记录。我还在追逐着我的小白兔。
